% =======================================================================
% Onward! 2027 — technical core of the language (draft, 2026-09-25)
%
% Written in response to the Onward! 2026 metareview: "The Trenza
% language ... is not sufficiently presented or described."
%
% House rule for this section: every claim about what the verifier
% accepts or rejects is backed by a listing in paper/onward2027/listings/
% and by trenza-core/tests/paper_listings.rs. Listings are included with
% \lstinputlisting from those very files, so the paper cannot show a
% listing that the compiler has not checked.
%
% Rule semantics below describe the validator at commit 3f5a336
% (trenza-core/src/validator.rs). Where the implementation is narrower
% than the 2026 paper claimed, this text follows the implementation; see
% history/chronicle/2026-09-25/01_CL_auditoria_afirmaciones_vs_codigo.md.
% =======================================================================

\section{The Language}
\label{sec:language}

This section presents Trenza concretely: a complete specification, the
grammar it is written in, the abstract model the verifier works on, and the
rules the verifier checks, stated as predicates over that model. Identifiers
in the listings are Spanish because they are taken from the case study
(Section~\ref{sec:validation}); \emph{Modo} means mode, \emph{tarjeta} card,
\emph{pesta\~na} tab.

\subsection{A Complete Specification}
\label{sec:example}

Listing~\ref{lst:modos} is a complete Trenza specification, accepted by the
verifier without diagnostics. It reduces the motivating defect of
Section~\ref{sec:motivation} to its essentials: a timer application with a
normal mode and an edit mode, in which the same three on-screen elements
react differently to the same tap.

\lstinputlisting[
  float=t,
  firstline=6,
  firstnumber=6,
  caption={A complete specification: two mutually exclusive modes.
           Verified without diagnostics
           (\texttt{paper/onward2027/listings/modos.trz}).},
  label={lst:modos}
]{onward2027/listings/modos.trz}

A specification has three layers. (Line numbers are those of the source
file, which begins with a comment header.) The \texttt{data} layer (lines 6--10)
declares what things \emph{are}: plain records with no behavior. The
\texttt{system} block (lines 12--17) declares the contexts, which one is
initial, and the external events the system receives. Each \texttt{context}
declares what things \emph{do there}: a \emph{role} binds a data type to a
name (\texttt{role tarjeta: Tarea}), and each handler
\texttt{on \emph{event} -> \emph{target}} states what that role does when it
receives that event in that context. A target is either an action call
(\texttt{iniciarTarea(self.tareaId)}), or one of two keywords:
\texttt{ignored} (the event is deliberately a no-op) and \texttt{forbidden}
(receiving the event is a runtime error). The \texttt{transitions} section
maps actions to the context that becomes active after them.

The point of the listing is what it does \emph{not} contain. There is no
\texttt{editMode} boolean and no conditional that tests it. Which behavior a
tap on a card has is determined by which context is active; the active
context is changed only by transitions. In the original JavaScript
implementation the same information was spread across several \texttt{if
(modoEdicion)} tests in different functions, and the defect was a place
that should have tested it and did not.
\ARGCESAR{Exact count: the 2026 paper says four places, the language spec
says five (app.js lines 286, 323, 411, 651, 662). Which one is right, and
was the defect the tab handler?}

\lstinputlisting[
  float=t,
  firstline=6,
  linerange={29-36},
  firstnumber=29,
  caption={The original defect, reintroduced: \texttt{ModoEdicion} no longer
           says what the frequent-items tab does on tap
           (\texttt{modos\_olvido.trz}). The verifier reports
           \texttt{[29:9] La acci\'on 'pestana\_frecuentes.tap' no est\'a
           declarada en el contexto 'ModoEdicion' [completeness]}
           and nothing else.},
  label={lst:olvido}
]{onward2027/listings/modos_olvido.trz}

Listing~\ref{lst:olvido} shows the same specification with the defect
reintroduced: the edit mode omits the handler for the frequent-items tab. In
the JavaScript original this omission was silent; the tab kept its
normal-mode behavior in edit mode. In Trenza it is rejected with exactly one
diagnostic, from Rule~1 below. Writing \texttt{on tap -> ignored} is not
boilerplate to silence the checker: it records that the author considered
the case and decided nothing should happen, and it makes that decision
visible to a reviewer, human or model.

\subsection{Grammar}
\label{sec:grammar}

Figure~\ref{fig:grammar} gives the core of the concrete syntax. It is a
simplified but faithful transcription of the PEG used by the compiler
(\texttt{trenza-core/src/trenza.pest}, 85 lines); the omitted productions
cover typed inputs, role bindings, decorators, lifecycle effects,
cross-package imports, and slots (a concurrent context contributing roles
to an overlay).
Layout is not significant; indentation in the listings is for the reader.

\begin{figure*}[t]
\small
\begin{align*}
\mathit{program}   &::= (\mathit{data} \mid \mathit{enum} \mid \mathit{system} \mid \mathit{context} \mid \mathit{external})^{*}\\
\mathit{data}      &::= \texttt{data}\ \mathit{Id}\ \mathit{annot}^{?}\ \texttt{:}\ (\texttt{mutable}^{?}\ \mathit{id}\ \texttt{:}\ \mathit{Type})^{*}\\
\mathit{system}    &::= \texttt{system}\ \mathit{Id}\ \texttt{:}\ \texttt{initial:}\ \mathit{Id}\ \mathit{section}^{*}\\
\mathit{section}   &::= \texttt{contexts:}\ \mathit{Id}^{*} \mid \texttt{overlays:}\ \mathit{Id}^{*} \mid \texttt{concurrent:}\ \mathit{Id}^{*} \mid \texttt{events:}\ \mathit{id}^{*}\\
\mathit{context}   &::= \texttt{context}\ \mathit{Id}\ \texttt{:}\ \mathit{clause}^{*}\\
\mathit{clause}    &::= \mathit{role} \mid \texttt{role * :}\ \mathit{target} \mid \texttt{transitions:}\ \mathit{trans}^{*} \mid \texttt{effects:}\ \mathit{effect}^{*} \mid \dots\\
\mathit{role}      &::= \texttt{role}\ \mathit{id}\ \mathit{annot}^{?}\ \texttt{:}\ \mathit{Type}\ (\texttt{on}\ \mathit{event}\ \texttt{->}\ \mathit{target})^{*}\\
\mathit{target}    &::= \texttt{ignored} \mid \texttt{forbidden} \mid \mathit{action}\ (\texttt{(}\ \mathit{args}\ \texttt{)})^{?}\\
\mathit{trans}     &::= \texttt{on}\ \mathit{action}\ \texttt{->}\ (\mathit{Id} \mid \texttt{[stay]} \mid \texttt{[close\_overlay]} \mid \texttt{[deactivate]})\\
\mathit{annot}     &::= \texttt{[}\ \mathit{key}\ \texttt{:}\ \mathit{value}\ \texttt{]}
\end{align*}
\caption{Core grammar of Trenza (simplified from the compiler's PEG).}
\label{fig:grammar}
\end{figure*}

\subsection{Abstract Model}
\label{sec:model}

The verifier does not reason about the generated code. It reasons about the
following abstraction of a specification, which the parser builds directly.
A specification denotes a tuple
\[
  S = (C,\ c_0,\ R,\ E,\ A,\ \tau,\ H,\ T,\ W)
\]
where $C$ is the finite set of contexts and $c_0 \in C$ the initial one;
$R$, $E$ and $A$ are the finite sets of role names, event names and action
names; $\tau : R \rightharpoonup \mathit{Type}$ assigns data types to roles;
\[
  H \subseteq C \times R \times E \times (A \cup \{\mathsf{ignored},\mathsf{forbidden}\})
\]
is the handler relation (one element per \texttt{on} line inside a role);
\[
  T \subseteq C \times A \times (C \cup \{\mathsf{stay},\mathsf{close},\mathsf{deactivate}\})
\]
is the transition relation; and $W \subseteq C$ is the set of contexts that
declare the wildcard \texttt{role *}. We write $\mathit{roles}(c)$ for the
roles declared in $c$ and
$\mathit{pairs}(c) = \{(r,e) \mid \exists t.\ (c,r,e,t) \in H\}$ for the
role--event pairs it handles.

The contexts named in the \texttt{system} block are classified as
\emph{base} (the initial context and those under \texttt{contexts:}; exactly
one is active), \emph{overlay} (stacked on top of the base, suspending it) or
\emph{concurrent} (active alongside the base). Every other context is a
\emph{sub-context}, and belongs to the overlay that opens it with
\texttt{initial:} or from which it is reached by transitions. This
classification induces a partition of $C$ into \emph{sibling groups}:
all base contexts form one group; the sub-contexts of each overlay form
one group; every overlay and concurrent context forms a group by itself.
We write $[c]$ for the group of $c$, and
$R_{[c]} = \bigcup_{c' \in [c]} \mathit{roles}(c')$,
$P_{[c]} = \bigcup_{c' \in [c]} \mathit{pairs}(c')$.

From $T$ the verifier derives the \emph{context graph}
$G = (C, \to)$ with $c \to c'$ iff $(c, a, c') \in T$ for some $a$, where
the pseudo-targets are abstracted as follows: $\mathsf{stay}$ adds no edge,
$\mathsf{close}$ and $\mathsf{deactivate}$ add an edge to $c_0$, and an
overlay with \texttt{initial: Sub} has an edge to \texttt{Sub}. The
\texttt{close} abstraction
abstraction is coarse (at run time, closing an overlay returns to whatever
context was beneath it, not necessarily $c_0$); Section~\ref{sec:limits}
discusses its consequences.

\subsection{The Rules}
\label{sec:rules}

Each rule is a predicate over $S$. A specification is accepted iff all
error-level predicates hold. All predicates are decidable in time
polynomial in the size of the specification; Rules 3 and 4 are graph
searches, the rest are set comparisons.

\begin{description}
\item[R1 Completeness.] A role--event pair handled in some context is
  handled in every sibling of that context that does not opt out:
  \[ \forall c \in C \setminus W.\ P_{[c]} \subseteq \mathit{pairs}(c). \]
\item[R2 Determinism.] Within a context, a role declares at most one
  handler per event: for every $(c,r,e)$ there is at most one \texttt{on}
  declaration, and hence at most one $t$ with $(c,r,e,t) \in H$.
\item[R3 Reachability] (warning). Every context is reachable in $G$:
  $\forall c \in C.\ c_0 \to^{*} c$.
\item[R4 Return.] From every context the initial one is reachable in $G$:
  $\forall c \in C.\ c \to^{*} c_0$. In particular R4 implies
  $G, c_0 \models \mathbf{AG}\,\mathbf{EF}\, c_0$ in CTL: along no path
  of $G$ does the system enter a context from which it cannot come back.
\item[R5 Role exhaustiveness.] A role that appears in some context
  appears in every sibling that does not opt out:
  $\forall c \in C \setminus W.\ \mathit{roles}(c) = R_{[c]}$.
\item[R6 Data conformance.] If a data type is annotated
  \texttt{[privacy: gdpr]}, a role may pass one of its fields as an action
  argument only if the role is annotated \texttt{[access: gdpr]}.
\item[R7 Slot integrity.] Every \texttt{fills X.s} names a slot \texttt{s}
  actually declared by context \texttt{X}; no slot is filled by two
  contexts; R2 holds inside each \texttt{fills} block.
\item[R8 Role-type consistency.] $\tau$ is a function: a role name has the
  same data type in every context in which it appears.
\end{description}

In addition, the verifier rejects user identifiers beginning with
\texttt{\_} (reserved for compiler-generated names), \texttt{initial:}
clauses that do not name a sub-context of an overlay, and imports whose
package does not define the imported system.

Two remarks keep these rules in proportion. First, R1, R2, R5 and R8 are
well-formedness conditions of the same nature as exhaustiveness checking of
pattern matches in ML-family languages; their value lies not in their
logical depth but in \emph{where} they are applied: to the table of
(context, role, event) behaviors that ordinary code spreads across
functions and files. Second, only R4 has a temporal reading, and only over
the abstraction $G$. We do not claim that the rules are equivalent to
arbitrary TLA$^{+}$ or Alloy specifications; they are a fixed, small set of
properties chosen because each one closes a defect class we observed.

\subsection{What Is Generated}
\label{sec:generated}

From an accepted specification the compiler emits: a Rust module (an enum
of contexts, a handler function per role--event pair, and a
\texttt{System} type that dispatches events and maintains the active
context, an overlay stack and a set of active concurrent contexts); a
TypeScript module with the same interface; one example-based test per
transition and per (context, role, event) handler, in Rust and in Vitest;
a Mermaid state diagram; and a Markdown report listing contexts,
transitions and the role--event table. For Listing~\ref{lst:modos}
\TODO{insert line counts and the generated \texttt{handle\_tarjeta\_tap}
function once the generator changes below are in}.

\subsection{Limits}
\label{sec:limits}

We state the limits of the current verifier explicitly, because each one
is a place where a specification can be accepted while the property a
reader would expect does not hold.

\paragraph{The wildcard.} A context that declares \texttt{role *: ignored}
is exempt from R1 and R5. The specification of Listing~\ref{lst:olvido}
with that single line added to \texttt{ModoEdicion} is accepted without
diagnostics (\texttt{modos\_comodin.trz}). An earlier version of the
verifier applied R1 and R5 to \emph{all} contexts, so that a modal dialog
had to declare a handler for every role of the main screen it covers; every
context of the case study then needed the wildcard, and R1 and R5 were
vacuous there. Scoping both rules to sibling groups removed the need for
it in 15 of the 18 contexts. The three that keep it are the phases of a
reset wizard, whose buttons differ from phase to phase; without the
wildcard they produce 24 diagnostics, which are design questions (should
a button absent from a phase be \texttt{ignored} or \texttt{forbidden}?)
rather than defects. \ARGCESAR{Decide the reset phases.}

\paragraph{Overlay return.} Abstracting $\mathsf{close}$ as an edge to
$c_0$ makes R4 sound only if every context beneath an overlay can itself
return to $c_0$, which R4 checks for base contexts but not for arbitrary
stacking. \TODO{State and check the stacking invariant.}

\paragraph{Dead transitions and indirect access.} R4 inspects declared
transitions but does not check that some handler of the context can produce
the triggering action, so a context whose only way back is never triggered
is still accepted. R6 recognises field accesses written as
\texttt{role.field} but not as \texttt{self.field}. Both gaps are pinned
by tests in \nolinkurl{trenza-core/tests/rules.rs}. \TODO{Close both.}

\paragraph{Generated code.} The verifier checks the specification, not the
generated code. The generated Rust currently ends every \texttt{match} with
a catch-all arm, so \texttt{rustc} does not independently re-check
completeness, and the Rust and TypeScript generators do not yet agree on
whether transitions are triggered by actions or by raw events.
\TODO{Remove the catch-all arms and align the two generators; then this
paragraph can be shortened.}
